\subsection{Base Allocation Heuristics}\label{subsec:heuristics}
\textit{(Joao, A)}\quad
Two resource-aware baseline heuristics were implemented for the integrated
simulation setting. Both methods receive the same filtered snapshot from the
simulator: resources are included only if they are currently available and not
busy, and tasks are included only if they are enabled, unassigned, unblocked and
compatible with the resource permissions. The heuristics are therefore reactive
with respect to availability. They do not predict future availability and are
included as transparent dispatching baselines, a common class of computationally
cheap resource-allocation rules in operational process support
contexts~\cite{park2023predictiveallocation,pufahl2025resourceallocation}.

\paragraph{Resource-aware Round Robin (R-RRA).}
R-RRA performs deterministic cyclic rotation over currently eligible and
available resources. It sorts resources by resource identifier and maintains a
persistent rotation pointer across allocation calls. From the rotated resource
order, each resource receives its oldest eligible waiting task, with task id used
as a deterministic tie-breaker. If a resource has no eligible task in the current
snapshot, it emits an idle decision and the rotation pointer advances only after
an actual assignment. Identical initial states and identical sequences of
allocation snapshots produce identical decision sequences.

\paragraph{Resource-aware Shortest Queue (R-SHQ).}
R-SHQ should be read as a cumulative-load-balancing heuristic rather than a
literal queue-length rule. The implementation uses
\texttt{ResourceEngine.load}, the cumulative number of assignments made to each
resource. This counter increments when a resource receives an assignment and is
not decremented when the resource is released; it therefore does not represent
the current queue length, remaining workload or future availability. Waiting
tasks are processed by increasing enabled time, then decreasing priority, then
task id. For each task, the method selects the feasible currently available
resource with the lowest cumulative assignment count and breaks ties by resource
id. Missing load values are treated as zero. The method is intentionally myopic:
it balances observed assignment counts at the moment of the decision, but it
does not optimize future queues, expected processing times or later arrivals.
Those consequences are evaluated only in Sec.~\ref{subsec:eval}.
